% TEX-root = ../main.tex
\chapter{\LaTeX{} 模板使用指南}

\section{环境配置}
\begin{enumerate}
    \item 本地环境\\
    下载安装最新版 TexLive 或 MikTeX。
    \item 在线环境\\
    Overleaf。
\end{enumerate}

\section{编辑文件}
\begin{enumerate}
    \item \texttt{sysusetup.tex}: 填写标题、作者、导师、学位名称等信息。
    \item \texttt{data/abstract.tex}: 填写中英文摘要。
    \item \texttt{data/denotation.tex}: 填写符号与缩略语，注意按音序排序。
    \item \texttt{data/chapxx.tex}: 各章内容，如有章节增删请在\texttt{main.tex}中修改相关记录。
    \item \texttt{data/appendix.tex}: 附录。
    \item \texttt{data/works.tex}: 学术成果。
    \item \texttt{data/acknowledgements.tex}: 致谢。
    \item \texttt{ref/refs.bib}: 引文数据库。
    \item \texttt{main.tex}: 主文件，用于控制文档选项（字体，学位类别）:
    \begin{enumerate}
        \item 学术硕士: \verb|\documentclass[degree=master]{sysuthesis}|;
        \item 专业硕士: \verb|\documentclass[degree=master,degree-type=professional]{sysuthesis}|;
        \item 博士: \verb|\documentclass[degree=doctor]{sysuthesis}|。
    \end{enumerate}
    指定论文要包括的部分，如摘要、目录、正文各章节、附录、引文数据库等等。
    （注意插入每章内容之后要加 \verb|\cleardoublepage| 以保证在打印版中各章都从右边开始）:
    \begin{verbatim}
    % !TeX root = ../main.tex

\chapter{学位论文形式结构}

\section{字数要求}
\begin{enumerate}
    \item 硕士论文: 正文一般为$1 \sim 3$万字；
    \item 博士论文: 正文一般不超过$15$万字。
\end{enumerate}


\section{前置部分}

\begin{enumerate}
    \item 封面和封底：由研究生院统一印刷，封面内容一律打印，其中导师姓名以研究生院备案名单为准。\\
          题目：在25字以内，能简明、具体、确切地表达论文特定内容，必要时，可以使用副标题。中文、英文题目应一致。
      \item 扉页内容包括学位论文中文、英文题目、专业名称、申请人姓名、导师姓名及论文答辩委员会组成 （由答辩委员会成员签名）。
      \item 原创性及学位论文使用授权声明。
      \item 中英文摘要。
\end{enumerate}


\section{主体部分}
\begin{enumerate}
    \item 主体部分包括引言 （前言），国内外文献综述，正文， 结语，参考文献。要求图表清晰，叙述流畅，章节有序，层次 分明。\\
    引言 （前言）部分内容主要为本研究课题的学术背景及理论与实际意义；本研究课题的来源及主要研究内容；建立研究的线索与思路。

    \item 文中的图、表、公式等，一律用阿拉伯数字按章顺序编号。如图 1-1、图2-2, 表 1-1、表 2-1，公式 (1-1) 等。图序及图名置于图的下方，居中排列；表序及表名置于表的上方，居中排列。详见第\ref{figures_tables}章的说明。

    \item 参考文献
    \begin{enumerate}
        \item 参考文献为论文中所有引文、引用观点以及对论文有重要影响和启发的文献；
        \item 参考文献按在论文中出现的先后依次排序；个别学科若通用该学科惯用的排序规范，可以例外；
        \item 参考文献内容一般排列在论文末尾 （论文篇幅较大且引用文献较多的，可在每章末尾注出），序码与论文加注处对应；
        \item 参考文献标注格式: 使用国标GB/T 7714-2015标准, 建议使用工具自动控制引文格式, 
        以保证格式规范。
        该\LaTeX{}模板已经对引文格式做了配置，
        用户需要将所需参考文献的信息存在 \texttt{bib} 格式的文件\texttt{ref/refs.bib}中, 
        通过 \verb|\cite{}|命令在恰当位置引用 (详见第\ref{citations}章的说明)。
        \texttt{bib}文件可从文献管理工具导出或自己用 \texttt{JabRef} 等软件编辑。
    \end{enumerate}

    \item 注释：可以用 “脚注”或 “文后注”来标注引用著作中的一些观点和案例，但全文标注方式应统一。
\end{enumerate}

\section{附录部分}
\begin{enumerate}
    \item 附录\\
        附录是正文主体的补充。下列内容可以作为附录：
        \begin{enumerate}
            \item 攻读学位期间发表的 （含已录用，并有录用通知书的）与学位论文相关的学术论文目录 。
            \item 由于篇幅过大，或取材于复制件不便编入正文的材料、数据。
            \item 对本专业同行有参考价值，但一般读者不必阅读的材料。
            \item 论文中使用的符号意义、单位缩写、程序全文及有关说明书。
            \item 附件：计算机程序清单、软磁盘、鉴定证书、获奖 奖状或专利证书的复印件等。
        \end{enumerate}
    \item 后记\\
        后记是有关本论文情况的说明性文字，主要是交代编写过程，阐述作者的感想和体会，对有关单位或个人的致谢语等。
\end{enumerate} 


    \cleardoublepage
    % !TeX root = ../thuthesis-example.tex

\chapter{图表示例}

\section{插图}

图片通常在 \env{figure} 环境中使用 \cs{includegraphics} 插入，如图~\ref{fig:example} 的源代码。
建议矢量图片使用 PDF 格式，比如数据可视化的绘图；
照片应使用 JPG 格式；
其他的栅格图应使用无损的 PNG 格式。
注意，LaTeX 不支持 TIFF 格式；EPS 格式已经过时。

\begin{figure}
  \centering
  \includegraphics[width=0.6\linewidth]{example-image-a.pdf}
  \caption{示例图片}
  \label{fig:example}
\end{figure}

若图或表中有附注，采用英文小写字母顺序编号，附注写在图或表的下方。
% LaTeX 传统上一般将附注的内容同图表的标题写在一起，形成很长的一段文字。

如果一个图由两个或两个以上分图组成时，各分图分别以(a)、(b)、(c)...... 作为图序，并须有分图题。
推荐使用 \pkg{subcaption} 宏包来处理， 比如图~\ref{fig:subfig-a} 和图~\ref{fig:subfig-b}。

\begin{figure}
  \centering
  \subcaptionbox{分图 A\label{fig:subfig-a}}
    {\includegraphics[width=0.45\linewidth]{example-image-a.pdf}}
  \subcaptionbox{分图 B\label{fig:subfig-b}}
    {\includegraphics[width=0.45\linewidth]{example-image-b.pdf}}
  \caption{多个分图的示例}
  \label{fig:multi-image}
\end{figure}



\section{表格}

表应具有自明性。为使表格简洁易读，尽可能采用三线表，如表~\ref{tab:three-line}。
三条线可以使用 \pkg{booktabs} 宏包提供的命令生成。

\begin{table}
  \centering
  \caption{三线表示例}
  \begin{tabular}{ll}
    \toprule
    文件名          & 描述                         \\
    \midrule
    thuthesis.dtx   & 模板的源文件，包括文档和注释 \\
    thuthesis.cls   & 模板文件                     \\
    thuthesis-*.bst & BibTeX 参考文献表样式文件    \\
    thuthesis-*.bbx & BibLaTeX 参考文献表样式文件  \\
    thuthesis-*.cbx & BibLaTeX 引用样式文件        \\
    \bottomrule
  \end{tabular}
  \label{tab:three-line}
\end{table}

表格如果有附注，尤其是需要在表格中进行标注时，可以使用 \pkg{threeparttable} 宏包。
研究生要求使用英文小写字母 a、b、c……顺序编号，本科生使用圈码 ①、②、③……编号。

\begin{table}
  \centering
  \begin{threeparttable}[c]
    \caption{带附注的表格示例}
    \label{tab:three-part-table}
    \begin{tabular}{ll}
      \toprule
      文件名                 & 描述                         \\
      \midrule
      thuthesis.dtx\tnote{a} & 模板的源文件，包括文档和注释 \\
      thuthesis.cls\tnote{b} & 模板文件                     \\
      thuthesis-*.bst        & BibTeX 参考文献表样式文件    \\
      thuthesis-*.bbx        & BibLaTeX 参考文献表样式文件  \\
      thuthesis-*.cbx        & BibLaTeX 引用样式文件        \\
      \bottomrule
    \end{tabular}
    \begin{tablenotes}
      \item [a] 可以通过 xelatex 编译生成模板的使用说明文档；
        使用 xetex 编译 \file{thuthesis.ins} 时则会从 \file{.dtx} 中去除掉文档和注释，得到精简的 \file{.cls} 文件。
      \item [b] 更新模板时，一定要记得编译生成 \file{.cls} 文件，否则编译论文时载入的依然是旧版的模板。
    \end{tablenotes}
  \end{threeparttable}
\end{table}

\cleardoublepage

    \cleardoublepage
    % !TeX root = ../main.tex

\chapter{引用文献的标注}\label{citations}

引用文献采用按照在文中被引用的顺序排序。
下文介绍 BibTeX 配合 \pkg{natbib} 宏包的主要使用方法。

在顺序编码制下，默认的 \cs{cite} 命令同 \cs{citep} 一样，序号置于方括号中，
引文页码会放在括号外。
统一处引用的连续序号会自动用短横线连接。

\sysusetup{
  cite-style = super,
}
\begin{tabular}{l@{\quad$\Rightarrow$\quad}l}
  \verb|\cite{zhangkun1994}|               & \cite{zhangkun1994}               \\
  \verb|\citet{zhangkun1994}|              & \citet{zhangkun1994}              \\
  \verb|\citep{zhangkun1994}|              & \citep{zhangkun1994}              \\
  \verb|\cite[42]{zhangkun1994}|           & \cite[42]{zhangkun1994}           \\
  \verb|\cite{zhangkun1994,zhukezhen1973}| & \cite{zhangkun1994,zhukezhen1973} \\
\end{tabular}


也可以取消上标格式，将数字序号作为文字的一部分。
建议全文统一使用相同的格式。

\begin{verbatim}
\sysusetup{
  cite-style = inline,
}
\end{verbatim}

\sysusetup{
  cite-style = inline,
}
\begin{tabular}{l@{\quad$\Rightarrow$\quad}l}
  \verb|\cite{zhangkun1994}|               & \cite{zhangkun1994}               \\
  \verb|\citet{zhangkun1994}|              & \citet{zhangkun1994}              \\
  \verb|\citep{zhangkun1994}|              & \citep{zhangkun1994}              \\
  \verb|\cite[42]{zhangkun1994}|           & \cite[42]{zhangkun1994}           \\
  \verb|\cite{zhangkun1994,zhukezhen1973}| & \cite{zhangkun1994,zhukezhen1973} \\
\end{tabular}



% \section{著者-出版年制}
% 
% 著者-出版年制下的 \cs{cite} 跟 \cs{citet} 一样。
% 
% \sysusetup{
%   cite-style = author-year,
% }
% \begin{tabular}{l@{\quad$\Rightarrow$\quad}l}
%   \verb|\cite{zhangkun1994}|                & \cite{zhangkun1994}                \\
%   \verb|\citet{zhangkun1994}|               & \citet{zhangkun1994}               \\
%   \verb|\citep{zhangkun1994}|               & \citep{zhangkun1994}               \\
%   \verb|\cite[42]{zhangkun1994}|            & \cite[42]{zhangkun1994}            \\
%   \verb|\citep{zhangkun1994,zhukezhen1973}| & \citep{zhangkun1994,zhukezhen1973} \\
% \end{tabular}

\vskip 2ex
\sysusetup{
  cite-style = super,
}
注意，引文参考文献的每条都要在正文中标注
\cite{zhangkun1994,zhukezhen1973,dupont1974bone,zhengkaiqing1987,%
  jiangxizhou1980,jianduju1994,merkt1995rotational,mellinger1996laser,%
  bixon1996dynamics,mahui1995,carlson1981two,taylor1983scanning,%
  taylor1981study,shimizu1983laser,atkinson1982experimental,%
  kusch1975perturbations,guangxi1993,huosini1989guwu,wangfuzhi1865songlun,%
  zhaoyaodong1998xinshidai,biaozhunhua2002tushu,chubanzhuanye2004,%
  who1970factors,peebles2001probability,baishunong1998zhiwu,%
  weinstein1974pathogenic,hanjiren1985lun,dizhi1936dizhi,%
  tushuguan1957tushuguanxue,aaas1883science,fugang2000fengsha,%
  xiaoyu2001chubanye,oclc2000about,scitor2000project%
}。

\cleardoublepage

    \cleardoublepage
    \end{verbatim}
\end{enumerate}

除上述文件外，
如无必要请勿修改其他重要文件，
如 \verb|sysuthesis-numeric.bst|（用于控制引文格式），
\verb|sysuthesis.cls| （文档类，用于控制文档显示的样式）等。

\section{编译文档}

可直接在命令行中使用 \texttt{latexmk} 命令，也可自己在编辑器中设置快捷键等。

为保证字体严格符合学校规定，建议终版文档在 Windows 平台上编译或在其他平台上安装 Windows 字体并修改\texttt{main.tex}中的文档选项，指定使用 Windows 字体。
\begin{verbatim}
    \documentclass[degree=doctor, fontset=windows]{sysuthesis}
\end{verbatim}